% =============================================================================
% 4.1 Experiment 1 Results — VM Degradation Profile (SQ1)
% =============================================================================

\section{Experiment 1 Results --- VM Degradation Profile (SQ1)}
\label{sec:results-exp1}

Experiment~1 ran the two-phase \texttt{thesis\_workload} job (30\,s CPU burn + 30\,s memory pressure with 400\,MB allocation, $\approx$60\,s total under no contention) on the Multipass VM (4\,vCPU, 4\,GB RAM) under the DockerRunLauncher across all six concurrency levels, three repetitions each.
Table~\ref{tab:results-vm} summarises the measured outcomes.

\begin{table}[H]
  \centering
  \caption{VM deployment results by workload level (3 repetitions, 4\,vCPU / 4\,GB RAM). Mean time and std are for \emph{successful} runs only. CPU and memory are system-wide averages during the run window.}
  \label{tab:results-vm}
  \small
  \begin{tabular}{clccccc}
    \toprule
    \textbf{Level} & \textbf{Jobs} & \textbf{SR\,(\%)} & \textbf{Mean\,(s)} & \textbf{Std\,(s)} & \textbf{CPU\,(\%)} & \textbf{Mem\,(\%)} \\
    \midrule
    L1  &  1 & 100.0 &  65.3 &  1.0 & 48.8 & 75.0 \\
    L2  &  2 & 100.0 &  76.2 & 13.4 & 60.8 & 81.6 \\
    L3  &  3 &  66.7 &  69.9 &  1.6 & 50.8 & 81.5 \\
    L5  &  5 &  46.7 &  76.3 &  3.2 & 46.7 & 81.4 \\
    L7  &  7 &  42.9 &  77.7 &  8.3 & 42.7 & 79.0 \\
    L10 & 10 &  30.0 &  79.9 &  1.4 & 38.0 & 77.3 \\
    \bottomrule
  \end{tabular}
\end{table}

\noindent\textit{SR = success rate across all launched runs (launched = $n \times 3$ per level). Mean time and Std are computed from successful runs only. CPU and Mem are system-wide averages from \texttt{vm\_metrics.csv} sampled every 1\,s during execution.}

\begin{figure}[H]
  \centering
  \includegraphics[width=0.85\textwidth]{vm_degradation.png}
  \caption{VM success rate and mean execution time vs.\ concurrency level.
  Success rate collapses from 100\,\% at L1--L2 to 30\,\% at L10, driven by OOM kills during the memory pressure phase.
  Execution time of \emph{successful} jobs remains comparatively stable (65--80\,s).}
  \label{fig:vm-degradation}
\end{figure}

\subsection{Observed Degradation Pattern}
\label{subsec:exp1-degradation}

Two distinct phases are visible in the results:

\paragraph{Phase 1 --- Reliable operation (L1--L2).}
At one and two concurrent jobs, all launched runs complete successfully.
Mean execution time is 65.3\,s (L1) and 76.2\,s (L2) --- a 17\,\% increase already attributable to initial CPU contention on the 4\,vCPU VM.
L2 exhibits notably high standard deviation (13.4\,s), indicating that even two concurrent jobs begin to produce variance in execution order and resource allocation timing.

\paragraph{Phase 2 --- OOM-driven failure (L3--L10).}
Beginning at three concurrent jobs, a fraction of jobs fail with short durations (37--50\,s) --- midway through the \texttt{memory\_pressure} op.
Inspection of individual run records reveals that failing jobs terminate at 37--50\,s, consistent with an OOM kill occurring within 7--20\,s of the memory allocation call in \texttt{memory\_pressure}.
The failure mechanism is \textbf{memory exhaustion}: with $n$ containers each allocating 400\,MB simultaneously, the aggregate demand quickly approaches or exceeds the 4\,GB VM RAM, triggering the Linux OOM killer which terminates one or more containers.

The success rate drops sharply from 100\,\% to 66.7\,\% between L2 and L3, and continues declining:
\begin{itemize}
  \item L3 (3 jobs): 6/9 = 66.7\,\%
  \item L5 (5 jobs): 7/15 = 46.7\,\%
  \item L7 (7 jobs): 9/21 = 42.9\,\%
  \item L10 (10 jobs): 9/30 = 30.0\,\%
\end{itemize}

\paragraph{Execution time of surviving jobs.}
The mean execution time of \emph{successful} runs stays comparatively flat across levels (65--80\,s).
This is a selection effect: only jobs that receive sufficient memory and CPU complete; these are typically the first-scheduled containers that acquire memory before the OOM killer is triggered.
The standard deviation among successful jobs ranges from 1.0\,s (L1) to 13.4\,s (L2), with moderate variance at L7 (8.3\,s).

\paragraph{CPU and memory trends.}
System CPU utilisation peaks at L2 (60.8\,\%) and then \emph{decreases} at higher levels (38\,\% at L10).
This counter-intuitive pattern reflects that at high concurrency the majority of jobs are OOM-killed within 40--50\,s, before the CPU-intensive workload contributes to sustained utilisation.
System memory remains elevated throughout (75--82\,\%) due to the combined footprint of the Docker daemon, PostgreSQL, Dagster daemon, and the run containers.

\paragraph{Reliability crossover reached.}
The reliability crossover (VM success rate falling below 95\,\%) occurs between \textbf{L2 and L3}: success rate drops from 100\,\% to 66.7\,\% as the aggregate memory demand first exceeds the VM's available RAM.
This is notably earlier than predicted by the hypothesis, which expected the crossover in the L5--L7 range.


\subsection{Resource Utilisation}
\label{subsec:exp1-resources}

System CPU utilisation peaks at L2 (60.8\,\%) and then decreases at higher levels (38.0\,\% at L10), because at high concurrency the majority of jobs are OOM-killed before completing the CPU-intensive workload.
Memory utilisation follows a similar pattern: each job allocates 400\,MB during its memory-pressure phase, so peak aggregate memory scales as $N \times 400\,\text{MB}$ per level.
At L10, aggregate memory allocation equals the VM's total 4\,GB RAM, explaining the sharp increase in failure rate at that level.

The key SQ1 finding is that the VM exhibited a \textbf{sharp reliability collapse at L3}: success rate dropped from 100\,\% to 66.7\% as aggregate memory demand first exceeded available RAM, and continued declining to 30\,\% at L10. This OOM-driven failure is the primary motivation for evaluating Kubernetes pod isolation.

